% Introdução (seção escrita pelos três).
% As hipóteses entram logo depois, de hipoteses_rq3_rq5.tex.

As métricas DORA, propostas pelo programa \textit{DevOps Research and
Assessment}~\cite{dora2024}, descrevem o desempenho de entrega de uma equipe por quatro
indicadores: frequência de \textit{deploy}, \textit{lead time} para mudanças, taxa de falha
das mudanças (CFR) e tempo de recuperação. Com os quatro, a equipe é classificada como
\textit{Elite}, \textit{High}, \textit{Medium} ou \textit{Low}, e a literatura associa as
categorias superiores a práticas como integração contínua e implantação em pequenos
lotes~\cite{forsgren2018}. As faixas vieram de questionários autorrelatados por equipes
corporativas, que não podem ser auditados, e não se sabe o que significam em projetos
\textit{open-source}, cujo ritmo de publicação é decidido por mantenedores voluntários.

Neste trabalho calculamos as quatro métricas apenas com dados públicos do GitHub, para 350
repositórios populares que usam GitHub Actions, em uma janela fixa de 12 meses. Como a API
não informa quando uma mudança chegou à produção nem quando um defeito afetou usuários,
usamos dois \textit{proxies}: a publicação de uma \textit{release} aproxima o
\textit{deploy}, e a falha de um \textit{workflow} de CI aproxima a falha de uma mudança.
Também verificamos esses \textit{proxies}: três avaliadores rotularam uma amostra-ouro de 60
repositórios e 299 \textit{releases}, contra a qual avaliamos a heurística que identifica
\textit{releases} corretivas, e comparamos três definições de \textit{deploy} para ver o
quanto as conclusões mudam com elas.

\subsection{Questões de pesquisa}

\begin{description}\itemsep0pt
\item[RQ 01] Com que frequência são feitos \textit{deploys} nesses projetos?
\item[RQ 02] Qual o \textit{lead time} para mudanças, nas variantes (a) e (b)?
\item[RQ 03] Qual a taxa de falha das mudanças, nas variantes (a) e (b)?
\item[RQ 04] Qual o tempo de recuperação após uma falha?
\item[RQ 05] Frequência de \textit{deploy} e taxa de falha se correlacionam?
\item[RQ 06] Que características dos repositórios se associam ao desempenho?
\item[RQ 07] A classificação muda se o \textit{deploy} for definido de outra forma?
\item[RQ 08] Qual a proporção de \textit{releases} corretivas (\textit{rework rate})?
\end{description}

A Seção~2 descreve a metodologia, a Seção~3 os resultados, a Seção~4 a discussão, a Seção~5
as ameaças à validade~\cite{wohlin2012} e a Seção~6 as conclusões.
